\section{Theoretical Analysis}
\label{appendix:theoretical-analysis}

This section presents a theoretical analysis of the memory capacity of \ours. We demonstrate that, given a set of queries exhibiting some degree of semantic similarity, our memorization rules in \Cref{eq:update-function} enable edge embeddings to memorize these queries and store sufficient information, thereby enabling effective memory replay. Formally, for any given edge, when the maximum angle $\theta$ between all query embedding pairs within a query set satisfies \Cref{eq:lamba-theta2}, the edge can effectively memorize the semantic information within that query set. 
\begin{equation}
    \theta \leq \lim_{d \to \infty} \left[ 2\arcsin\left(\sqrt{\frac{1}{2}}\sin(\arccos(\lambda))\right) \right]
\label{eq:lamba-theta2}
\end{equation}


\begin{assumption}
For a typical path $e = A \rightarrow B$ in the knowledge graph where $e \in E$, we define that this path "remembers\footnote{We ignore the first term in \Cref{eq:dfs-step-in}, as it mainly serves to adjust the system's sensitivity to memorized data and is multiplied by a small coefficient, rendering its overall impact on the final result relatively minor.}" query $q$ when $q \cdot e > \lambda$.
We assume that the input query embedding $q$ is normalized, i.e., $\|q\| = 1$.
\end{assumption}

\begin{proposition}
For a set of query embeddings $Q\subseteq \mathbb{R}^n$, we have
\begin{align*}
&\forall q_1, q_2 \in Q, \angle(q_1, q_2) \leq 2\arcsin\left(\sqrt{\frac{1}{2}}\sin(\arccos(\lambda))\right) \implies \exists n \in \mathbb{N}, \exists (q_1, q_2, \dots, q_n) \in Q^n, \\
&\quad\quad\quad\quad \text{such that }\quad e_{\text{final}} = (f_{q_n} \circ \dots \circ f_{q_1})(\vec{0}) 
\text{ satisfies } \forall q \in Q,\quad e_{\text{final}} \cdot q > \lambda
\end{align*}
\end{proposition}


\begin{proof}
\label{pr}
Each query embedding $q \in Q$ can be represented as a point on the $d$-dimensional unit sphere $\mathbb{S}^{d-1}$, i.e., there exists a bijective mapping $\varphi : Q \to N \subseteq \mathbb{S}^{d-1}(0,1)$ for all $q \in Q$.

Assuming $\forall q_1, q_2 \in Q$, their angle satisfies $\angle(q_1, q_2) < \theta$, the corresponding point set $N$ satisfies $\forall p_1, p_2 \in N, d_{\mathbb{S}^{d-1}}(p_1, p_2) < \theta$, where $d_{\mathbb{S}^{d-1}}$ denotes the geodesic distance on $\mathbb{S}^{d-1}$.

According to \textbf{Jung's Theorem on Spheres} \cite{jung-theory1,jung-theory2} in $\mathbb{S}^{d-1}$, there exists a ball $P = \mathbb{B}^{d-1}(c_p,r_p)$ containing $N$ whose geodesic radius $r_P$ satisfies:
\[
\theta \geq 2\arcsin\left(\sqrt{\frac{d+1}{2d}}\sin(r_P)\right) 
\quad \xrightarrow{d \to \infty} \quad 
\theta \geq 2\arcsin\left(\sqrt{\frac{1}{2}}\sin(r_P)\right).
\]
Since we typically deal with very high embedding dimensions in practice, we can approximate $d \to \infty$.

From Jung's Theorem, the center $c_p$ of the covering ball $P$ lies within the convex hull of $N$. Therefore, the unit vector $c = \varphi^{-1}(c_p)$ satisfies:
\begin{enumerate}[label=\arabic*)]
\item $\forall q \in Q, \angle(c, q) \leq r_P \implies \langle c, q \rangle \geq \cos(r_P).$
\item There exist coefficients $\{\eta_i\}_{i=1}^k \subseteq \mathbb{R}_+$ with $\sum_{i=1}^k \eta_i = 1$ such that $c = \sum_{i=1}^k \eta_i q_i.$
\end{enumerate}

From (2), it follows that there exists an update sequence $\{q_{i_j}\}_{j=1}^m \subseteq Q$ such that:
\[
e_{\text{final}} = \lim_{m \to \infty} (f_{q_{i_m}} \circ \cdots \circ f_{q_{i_1}})(\mathbf{0}) = c.
\]
Since $f(x) = 2\arcsin\left(\sqrt{\frac{1}{2}}\sin(r_P)\right)$ is monotonically decreasing:
\[
\lambda \leq \cos(r_P) \implies 2\arcsin\left(\sqrt{\frac{1}{2}}\sin(\arccos(\lambda))\right) \geq \theta \implies \forall q \in Q, \langle e_{\text{final}}, q \rangle \geq \lambda,
\]
which satisfies the proposition's requirements.
\end{proof}